\capitulo{2}{Objetivos del proyecto}

El presente Trabajo de Fin de Grado tiene su origen en una propuesta planteada
por el equipo de tutorización académica de la Escuela Politécnica Superior,
aceptada por el alumno motivado por su interés particular en profundizar en el
estudio de las redes neuronales y los algoritmos de \textit{Machine Learning}.
La aplicación de estas técnicas a un caso de uso real, como es el análisis de
muestras de miel mediante espectrometría, ha permitido conjugar un objetivo
formativo de especialización técnica con la resolución de una problemática
concreta del sector apícola, en el cual los análisis convencionales suelen
presentar costes elevados y tiempos de procesamiento prolongados.

Conforme a la estructura habitual de proyectos de Ingeniería del Software, los
objetivos planteados en este TFG se han clasificado en dos categorías
diferenciadas: por un lado, los \textbf{objetivos del producto software}, que
definen el comportamiento esperado del sistema resultante; y por otro, los
\textbf{objetivos técnicos}, que recogen los retos formativos y de
implementación abordados durante el desarrollo del proyecto.

\section{Objetivos del producto software}

Los objetivos del producto definen el alcance funcional del sistema HIVES
(\textit{Honey Identification and Verification by Spectroscopy}), cuya
especificación detallada se encuentra en el Apéndice B. A continuación,
se presenta una síntesis de los mismos:

\begin{description}
    \item[OP-1: Automatización del proceso de análisis.] \leavevmode\\ Habilitar la captura
    automatizada de los datos espectrales mediante la integración del software
    con un dispositivo espectrométrico, eliminando la necesidad de intervención
    manual durante la lectura de las muestras.

    \item[OP-2: Clasificación inteligente mediante técnicas de IA.] \leavevmode\\ Implementar
    un sistema de inferencia basado en una red neuronal preentrenada con
    \textbf{TensorFlow} o \textbf{Scikit-Learn}, capaz de discriminar el tipo de miel a partir de su
    firma espectral con un nivel de confianza cuantificable.

    \item[OP-3: Generación automatizada de informes técnicos.] \leavevmode\\ Producir reportes
    en formato \textbf{PDF} que recojan la totalidad de la información relevante
    de cada análisis (fecha, resultados de clasificación, métricas espectrales y
    métricas de confianza del modelo), facilitando su archivo y distribución.

    \item[OP-4: Persistencia y trazabilidad del histórico.] \leavevmode\\ Almacenar los datos
    de cada análisis en una base de datos local, garantizando su consulta
    posterior y posibilitando la extracción de conjuntos de datos para futuros
    reentrenamientos del modelo.

    \item[OP-5: Usabilidad y accesibilidad operativa.] \leavevmode\\ Proporcionar una interfaz
    gráfica intuitiva que permita el manejo del sistema por parte
    de personal sin conocimientos técnicos avanzados, asegurando además la
    compatibilidad multiplataforma.
\end{description}

\section{Objetivos técnicos}

Los objetivos técnicos recogen los principales retos formativos, metodológicos
y de implementación que se han afrontado durante el desarrollo del proyecto.
Estos objetivos no se traducen directamente en funcionalidades del producto,
sino que constituyen el ámbito de aprendizaje y especialización del alumno como
ingeniero informático.

\begin{description}
    \item[{\parbox[t]{\linewidth}{OT-1: Profundizar en el estudio de algoritmos de \textit{Machine Learning}.}}] \leavevmode\\\leavevmode\\
    Adquirir un conocimiento sólido sobre los
    fundamentos teóricos y prácticos del aprendizaje automático, abordando la
    totalidad del ciclo de vida del modelo: preparación y normalización de los
    datos, diseño de la arquitectura de la red, entrenamiento, evaluación de
    métricas e integración del modelo entrenado en una aplicación de producción.

    \item[{\parbox[t]{\linewidth}{OT-2: Emplear \textit{Jupyter Notebooks} como entorno de experimentación.}}]  \leavevmode\\\leavevmode\\
    Utilizar el ecosistema de cuadernos interactivos para
    llevar a cabo las fases de exploración de los datos espectrales, prototipado
    de arquitecturas de red, entrenamiento iterativo y evaluación comparativa
    de los modelos, separando de forma clara la fase de investigación de la
    posterior fase de integración en la aplicación final.

    \item[{\parbox[t]{\linewidth}{OT-3: Integrar la comunicación hardware-software a bajo nivel.}}] \leavevmode\\\leavevmode\\
    Establecer una interconexión fiable entre la aplicación de escritorio y el
    dispositivo espectrométrico mediante el protocolo de comunicación por
    puerto serie, gestionando la sincronización, los tiempos de espera y los
    escenarios de desconexión abrupta del hardware.

    \item[{\parbox[t]{\linewidth}{OT-4: Aplicar el patrón arquitectónico Modelo-Vista-Controlador (MVC).}}] \leavevmode\\\leavevmode\\
    Diseñar e implementar el sistema siguiendo una arquitectura modular
    que garantice una separación estricta entre la interfaz gráfica, la lógica
    de negocio y la capa de datos, favoreciendo la mantenibilidad y la
    escalabilidad del software.

    \item[{\parbox[t]{\linewidth}{OT-5: Gestionar la concurrencia en aplicaciones de escritorio.}}] \leavevmode\\\leavevmode\\
    Implementar un modelo de ejecución multihilo que permita delegar las tareas
    de adquisición de datos y de inferencia del modelo a hilos secundarios,
    evitando el bloqueo de la interfaz gráfica (\textit{GUI freezing}) durante
    las operaciones intensivas.

    \item[{\parbox[t]{\linewidth}{OT-6: Aplicar una metodología ágil real en un proyecto unipersonal.}}] \leavevmode\\\leavevmode\\
    Llevar a la práctica el marco de trabajo \textbf{Scrum} apoyado en la
    herramienta \textbf{Jira}, asumiendo la gestión completa del
    \textit{Product Backlog}, la planificación de \textit{sprints}, la
    estimación mediante \textit{Story Points} y el análisis de las métricas de
    velocidad y \textit{burn-down}.

    \item[{\parbox[t]{\linewidth}{OT-7: Garantizar la compatibilidad multiplataforma.}}] \leavevmode\\ \leavevmode\\
    Diseñar la aplicación de forma que pueda ejecutarse de manera consistente sobre los
    sistemas operativos Windows, Linux y macOS, gestionando adecuadamente las
    particularidades de cada entorno en lo relativo al acceso al puerto serie
    y al sistema de archivos.

    \item[{\parbox[t]{\linewidth}{OT-8: Diseñar una capa de persistencia local eficiente.}}] \leavevmode\\\leavevmode\\
    Implementar el almacenamiento de los datos mediante el motor \textbf{SQLite},
    prescindiendo del uso de mapeadores objeto-relacional (ORM) con el fin de
    optimizar el rendimiento y mantener un control directo sobre las sentencias
    SQL.

    \item[{\parbox[t]{\linewidth}{OT-9: Empaquetar y distribuir la aplicación como ejecutable independiente.}}] \leavevmode\\\leavevmode\\ 
    Generar un binario autónomo que encapsule el intérprete de
    Python, las dependencias y el modelo entrenado, permitiendo al usuario
    final ejecutar el sistema sin requerir instalaciones técnicas previas.
\end{description}

La consecución conjunta de los objetivos del producto y los objetivos técnicos
descritos constituye la base sobre la que se ha articulado el diseño, la
implementación y la validación del sistema HIVES, cuyos detalles se desarrollan
en los apéndices posteriores.
